\section{Introduction}
\label{sec:intro}

Solvers for mixed-integer nonlinear programs repeatedly solve convex
subproblems and act on the answers. The optimal value of a relaxation
decides whether a node of the search tree can be pruned. The location of an
optimizer guides branching and rounding, and whether a bound or constraint is
active at that optimizer informs the choice of branching variables. A
certificate of optimality is what an independent checker verifies. In
floating-point arithmetic all of these quantities are approximations. Exact
and proof-producing solvers need exact versions of them. For linear and
convex quadratic subproblems the exact versions are available: a convex
quadratic program has a rational optimizer of polynomial bit length, and it
can be computed in polynomial time~\cite{KozlovTarasovKhachiyan1980}. From
degree four on, a globally convex objective can have a unique optimizer that
is irrational, and the different outputs come apart.

This paper studies polynomial optimization problems with explicitly encoded
rational data. Most results assume convexity of one of three kinds, stated in
each theorem: global convexity or strong convexity of the objective,
convexity on the feasible set only, or, in Section~\ref{sec:recourse},
convexity in a subset of the variables. Appendix~\ref{app:further-arithmetic}
also treats nonconvex quadratic constraints and cone constraints with
explicit algebraic data. We ask which of the following
outputs can be computed or represented efficiently, and what makes the
others hard:
\begin{enumerate}[label=(\roman*)]
\item the optimal value to a requested precision $2^{-q}$;
\item an optimizer to precision $2^{-q}$, preferably the same optimizer at
  every precision;
\item an exact decision about the optimizer, such as the sign of a given
  polynomial at the optimizer, the comparison of the optimal value with a
  rational threshold, or whether a constraint is active;
\item an exact description of the optimizer, or of a certificate of its
  optimality.
\end{enumerate}
Section~\ref{sec:models} defines these outputs, the input encodings, and the
computational models. Throughout, $L$ denotes the total binary length of the
input, $n$ the number of continuous variables, $k$ a structural parameter
specified in each statement, $q$ the number of requested bits of precision,
$f$ the objective, and $p$ the selected optimizer, which is the unique
optimizer when the objective is strongly convex and the feasible set is
convex.

Output~(i) is well understood for globally convex polynomials. Slot,
Steurer, and Wiedmer consider globally convex polynomials with rational
coefficients written in binary and exponent vectors written in unary. On a
nonempty rational polyhedron such a polynomial is either unbounded below or attains
its minimum at a point of norm at most
$2^{\poly(L)}$~\cite[Theorem~1.1]{SlotSteurerWiedmer2025}, and its optimal
value, when finite, can be approximated to objective gap $2^{-q}$ in time
polynomial in the input length and $q$~\cite[Corollary~1.2]{SlotSteurerWiedmer2025}. The norm bound does not
bound the bit length of the coordinates of an optimizer, which can be
irrational. In arXiv version 2511.03440v1 of November~5, 2025, two questions
about convex quartics are marked open: whether every rational polyhedron that
contains a point with $f(x)\le0$ also contains such a rational point of
polynomial bit length, and the complexity of deciding exactly whether such a
point exists. All comparisons with that paper below refer to this version.
For outputs (iii) and~(iv) the arithmetic benchmark is the
problem $\PosSLP$~\cite{AllenderEtAl2009}: given a division-free arithmetic
circuit with constants $0$ and $1$, decide whether the integer it computes is
positive. Such an integer can have exponentially many bits. $\PosSLP$ lies in
the counting hierarchy~\cite[Theorem~1.4]{AllenderEtAl2009}, and the Square
Root Sum problem is decidable in polynomial time with a $\PosSLP$
oracle~\cite[Section~1.4]{AllenderEtAl2009}. Whether $\PosSLP$ is in P, or is
NP-hard, is open~\cite{BuergisserJindal2024}. Exact feasibility of
semidefinite programs is already at least as hard as
$\PosSLP$~\cite{TarasovVyalyi2008}.

The four outputs form a ladder: each step up needs an ingredient that the
step below does not supply. Our results identify these ingredients and the
classes of problems in which they are available or provably missing.
\begin{itemize}
\item \emph{From values to points.} An effective error bound is needed, that
  is, a bound on the distance to the optimal set in terms of the objective
  gap, with a constant computable from the input, together with a rule that
  selects one optimizer. Global convexity supplies such a bound at every
  degree, and convexity on a polytope supplies one for cubics. For quartics
  that are convex only on a box, a single constant-accuracy point already
  decides Square Root Sum or $\PosSLP$ (Section~\ref{sec:intro-points}).
\item \emph{From points to exact decisions.} A separation bound for nonzero
  values is needed, together with a compact way to reach the precision that
  it demands. For explicit globally strongly convex objectives of unary
  degree, every exact comparison at the minimizer reduces to one $\PosSLP$
  instance, and already for certified quartics the order comparisons are
  $\PosSLP$-hard (Section~\ref{sec:intro-exact}). Linear constraints add a
  search for the active face, and integer variables add a discrete search
  (Section~\ref{sec:intro-constraints}).
\item \emph{From exact decisions to exact descriptions.} Separate degree,
  height, and field arguments are needed, and the size of a description
  depends on its format. Expanded rational coordinates, coefficient lists of
  minimal polynomials, Gram matrices, and sum-of-squares certificates can be
  exponentially long in the dimension for certified inputs of polynomial
  size, while the exhibited families have short descriptions in compact
  formats such as circuits (Sections~\ref{sec:intro-optimizers}
  and~\ref{sec:intro-certificates}).
\end{itemize}
The central result is Theorem~\ref{thm:quartic-complete}: for explicit
quartics with a checkable strict convexity certificate, exact order
comparisons of the minimum value, and of a minimizer coordinate, with a
rational threshold are $\PosSLP$-complete. Together with
Theorem~\ref{thm:exact-upper}, it shows that for explicit globally strongly
convex polynomials of unary degree, exact order comparisons at the minimizer
are as hard as $\PosSLP$ and no harder; for equality only the upper bound is
known. Table~\ref{tab:intro-results} summarizes the results, and
Section~\ref{sec:intro-prior} separates them from prior work.

\subsection{From values to points}
\label{sec:intro-points}

Global convexity gives an effective conversion from values to points.
Theorem~\ref{thm:global-point} treats an explicitly sparse rational
polynomial $f$ that is convex on $\R^n$, of degree at most $D$, and an
arbitrary rational polyhedron $P$. In time polynomial in $L$ and $D$ one can
decide whether $P$ is empty and whether $f$ is unbounded below on $P$.
Otherwise the minimum $f_*$ is attained on a set $S$, and one can compute
rationals $R$ and $\Gamma$ of bit length $\poly(L,D)$ such that the
minimum-norm optimizer satisfies $\norm{p}\le R$ and
\begin{equation}
\label{eq:intro-error-bound}
\dist(x,S)\le\Gamma\,\bigl(f(x)-f_*\bigr)^{1/D}
\qquad(x\in P,\ f(x)-f_*\le1).
\end{equation}
Consequently this fixed optimizer can be approximated to precision $2^{-q}$
in time $\poly(L,D,q)$. Nonunique optima, singular Hessians, and unbounded
polyhedra are allowed. Global error bounds for convex polynomials are
known~\cite{Li2010,Li2013,Yang2009,Ngai2015}, and we do not claim the form of
the exponent as new. Our addition is a constant that is computable from the
input and has polynomial bit length on arbitrary rational polyhedra,
together with the fixed-selector output that it makes possible. With a
supplied strong convexity modulus the conversion from values to points is
immediate; the theorem matters when the objective has flat directions or
several minimizers.

Convexity on the feasible set alone suffices for cubics.
Theorem~\ref{thm:cubic-point} treats a cubic that is convex on a bounded
rational polytope, possibly of lower dimension, but not necessarily on
$\R^n$. It proves \eqref{eq:intro-error-bound} with exponent $1/4$ and a
constant of polynomial bit length, and approximates both the optimizer set
and its minimum-norm point in time $\poly(L+q)$. A cubic that is convex on all
of $\R^n$ is a quadratic, so the domain-convex case is the one that matters
for degree three. The univariate cubic $t^3$ on $[0,1]$ shows that no exponent
larger than $1/3$ is possible.

For quartics, convexity on the feasible set is not enough.
Section~\ref{sec:points} constructs rational quartics that are convex on a
box and have a unique minimizer, for which a single approximation of the
minimizer to accuracy $1/4$ decides the Square Root Sum problem or $\PosSLP$.
For Square Root Sum the constructed quartics have all coefficients other than
the constant term bounded by an absolute constant, and their interaction
graphs, which join two variables when some monomial contains both, have
treewidth at most two. A deterministic polynomial-time point algorithm for
this class would therefore put the corresponding problem in P.
Unconditionally, there are explicit box-convex quartics in $n+1$ variables,
easy to solve by other means, for which every error bound of the form
\eqref{eq:intro-error-bound} with exponent $\vartheta$ has a constant $\Gamma$
with $\log_2\Gamma\ge\vartheta\,2^{n+2}$, and for which isotropic
regularization approximates the optimizer to accuracy $1/2$ only with a
parameter whose denominator has more than $2^{n+2}$ bits. These instances
have input length $L=\Theta(n^2)$, so the two quantities are exponential in
the dimension and of order $2^{\Omega(\sqrt L)}$: superpolynomial, but not
exponential, in $L$. The same section separates finer output contracts. For a
strongly convex cubic on a box, deciding whether a designated coordinate of
the optimizer equals its lower bound is Square Root Sum-hard, although
Theorem~\ref{thm:cubic-point} approximates the optimizer in polynomial time.
There are also box-convex quartics for which some optimizer is easy to
approximate, while approximating the minimum-norm optimizer to constant
accuracy decides Square Root Sum.

The point problem also changes when part of the objective is nonconvex.
Section~\ref{sec:recourse} considers a cubic $F(v,z)$ on
$[0,1]^k\times[0,1]^m$ that is convex in the residual variables $z$ for each
value of the core variables $v$, but possibly nonconvex in $v$. Residual
values at rational cores are easy to compute, but approximations to an
optimal core do not determine the optimal residual: in the cubic
$v^2(2-z)$ on the unit square, the minimum-norm residual optimizer is $0$ at
$v=0$ and $1$ at every $v>0$. We perturb only the core, by a linear term
$\gamma^{\mathsf T}v$ drawn once from the uniform law on a finite rational
grid in $[-\sigma,\sigma]^k$. For every precision $q$, an evaluator returns
the lexicographically first optimal core of the perturbed objective together
with its minimum-norm residual optimizer, within $2^{-q}$, and every returned
point is correct, whatever the draw. The expected work is
$k^{O(k)}(1+\Lambda/\sigma)^k\poly(L+q)$, where
$\partial^2F/\partial v_i^2\le\Lambda$ along each core coordinate on the
product box, and one
random factor governs all precisions. Beyond a certified approximation of the
selected core, the proof certifies which core coordinates lie at their
bounds, by comparing slightly tilted problems, certifies margins of the
relevant minors by lattice basis reduction, and then completes the residual
by regularization. For every fixed degree, the same section also shows the
following. Certified values and the selected core are available under
residual convexity on a product box, or under a supplied convexifier on a
polytope that couples core and residual variables; a convexifier is a
rational $\alpha\ge0$ such that adding $\frac\alpha2\norm v^2$ makes the
objective convex there. A supplied convexifier also yields full points,
provided that the convexified objective is convex on the polytope and has
degree three, or is convex on all of $\R^{k+m}$. Quadratic objectives permit
exact rational optimizer and value recovery, by a polynomial height bound
and rational reconstruction (Corollary~\ref{cor:recourse-qp}). A positive
semidefinite quadratic plus positive even powers of rational affine forms
gives a checkable subclass of the global convexifier hypothesis
(Corollary~\ref{cor:recourse-affine-power}). With $\alpha=0$, that is,
for objectives that are jointly convex, the expected work in these results is
polynomial, with no factor $k^{O(k)}$ and no numerical dependence on $\sigma$.
For residual quartics no full-point guarantee of this form is expected:
such a guarantee, applied after adding a dummy core to the box-convex quartic
reductions above, would give a Las Vegas algorithm with expected polynomial
running time for Square Root Sum and for $\PosSLP$
(Remark~\ref{rem:recourse-quartic}).

\subsection{From points to exact decisions}
\label{sec:intro-exact}

Exact decisions require arithmetic beyond approximation.
Theorem~\ref{thm:exact-upper} treats explicit polynomials $f$ and $h$ of
unary degree together with a rational $\mu>0$ such that
$\nabla^2f\succeq\mu I$ on $\R^n$; the curvature bound is either a promise or
derived from a checked certificate. For each relation
${\bowtie}\in\{<,\le,=,\ne,\ge,>\}$, a deterministic polynomial-time algorithm,
which makes no oracle calls, constructs one integer arithmetic circuit whose
value is positive exactly when $h(p)\bowtie0$. No gap between $h(p)$ and zero
is assumed. The construction obtains a starting point of polynomial bit
length from convex value approximation, stores polynomially many exact
Newton steps as a shared rational circuit, and uses real quantifier
elimination to show that a nonzero value $h(p)$ has magnitude at least
$2^{-2^{e(L)}}$ for a fixed polynomial $e$. The Newton iterates reach this
precision because their error squares at each step. Shifted comparisons
handle the six relations, and division is removed by tracking numerators
and positive denominators. Section~\ref{sec:upper} proves the same reduction
for observables at the unique zero of a globally strongly monotone polynomial
map of unary degree whose Jacobian need not be symmetric. The gradients
of certified quartics are such cubic maps, so order tests are
$\PosSLP$-hard already for these gradient maps and complete for the
certified cubic-map class. Equality retains only the upper bound. Cubic maps are
the first nontrivial case, because a strongly monotone quadratic map is
affine.

The matching lower bound holds for explicitly certified quartics. A
\emph{full Hessian Gram} for a quartic $f$ is a rational
symmetric matrix $A$ with
$v^{\mathsf T}\nabla^2f(x)v=(v,x\otimes v)^{\mathsf T}A(v,x\otimes v)$; when
$A\succ0$ it certifies global strong convexity and can be verified in
polynomial time (Section~\ref{sec:models}). Theorem~\ref{thm:quartic-complete}
proves that, for explicit rational quartics supplied with a positive definite
full Hessian Gram, strict and weak comparisons of the minimum value with a
rational threshold, and of a coordinate of the minimizer with a rational
threshold, are $\PosSLP$-complete under polynomial-time many-one reductions.
For the value comparisons the hard instances have nonzero minima. In the same
class, nonnegativity, membership in the cone of real sums of squares, and
existence of a positive definite rational polynomial Gram matrix are
$\PosSLP$-complete, and membership in the cone of rational sums of squares is
$\PosSLP$-hard. Equality tests have the upper bound of
Theorem~\ref{thm:exact-upper} and no matching lower bound.
Section~\ref{sec:reductions} shows further that coordinate order comparisons
remain $\PosSLP$-complete when the minimizer is rational and lies in
$[-1,1]^n$, the minimum is known to be zero, and $n+1$ short rational square
factors are supplied. Rationality of the minimizer is then a promise of the
problem, not a property the reduction recognizes.

In particular, deciding whether $f(x)\le0$ has a real solution, that is,
whether $f_*\le0$, is $\PosSLP$-complete for certified quartics. The exact
feasibility question for convex quartics that is marked open in the arXiv
version of~\cite{SlotSteurerWiedmer2025} is therefore $\PosSLP$-hard, already
without constraints. Corollary~\ref{cor:reductions-feasibility} gives
the same hardness over a halfspace and over the unit cube, and for certified
quartics on arbitrary rational polyhedra Section~\ref{sec:constraints} gives
the upper bound $\mathrm{UP}^{\PosSLP}\cap\mathrm{coUP}^{\PosSLP}$. For convex
quartics that are not strongly convex, no upper bound relative to $\PosSLP$ is
known.

The main tool of the lower bounds is Lemma~\ref{lem:quartic-realization}.
Let rational quadratics $r_1,\ldots,r_m$ vanish at a point $a\in\R^n$ with
well-conditioned gradients there, and let a rational quadratic $G$ vanish at
$a$, have a positive definite quadratic part, and have a sufficiently small
gradient at $a$. Then
\[
F=\Bigl(\frac{G}{t\nu}\Bigr)^2+\sum_{j=1}^m\Bigl(\frac{r_j}{\nu}\Bigr)^2
\]
is a rational sum of $m+1$ squares, has unique zero $a$, satisfies
$\nabla^2F\succeq\frac32I$ on $\R^n$, and has a positive definite rational
full Hessian Gram, for suitable rationals $t$ and $\nu$ determined by
explicit quantitative bounds on the data. The Gram matrix is computed
directly from the rational square factors: canonical Hessian Gram matrices of
quadratic squares are exactly covariant under translation, so once the
factors are fixed, the certificate requires no further approximation of the
point $a$, which need not even be known. Squaring the residuals alone does not
give convexity: $x^2+x^4+\varepsilon(y-x^2)^2$ has second derivative $-2$ in
$x$ at $(0,1/\varepsilon)$. Section~\ref{sec:reductions} applies the lemma to
two simulations of integer arithmetic. The first uses cube roots of numbers
near one, in which an integer $V$ is encoded by the leading term $V\delta^d$
of a tiny analytic signal. The second uses products of rational unit
quaternions, which keeps every coordinate of the realized point rational.

Section~\ref{sec:upper} also proves a compilation result that several
constrained results use. Consider an integer circuit of size $S$ with Boolean
inputs, gates $+$, $-$, $\times$, and threshold gates $v\mapsto[v>0]$, whose
output is Boolean. In polynomial time it can be converted into a
division-free integer circuit of size $O(S^2)$, without threshold gates,
whose value is positive exactly when the original output is one. Each
threshold gate is replaced by a rational approximation of the sign function:
scaled maps $x\mapsto2\sigma x/(\sigma+x^2)$ first compress the range, and the
inverse Newton iteration $x\mapsto2x/(1+x^2)$ then refines the sign. The
exact gate values are integers of absolute value at most $2^{2^S}$, so
distinct values are at least one apart, and a single error budget keeps every
approximation within $1/4$ of its exact value; numerators with positive
denominators remove the divisions. A padded interpreter for query
descriptions extends the construction to every deterministic polynomial-time
computation with a $\PosSLP$ oracle, which therefore reduces to one $\PosSLP$
instance (Theorem~\ref{thm:posslp-closure}). We use this to state
deterministic oracle algorithms as single-instance reductions. The
construction compiles one output bit at a time and does not apply to
nondeterministic or randomized oracle algorithms.

\subsection{Constraints and integer variables}
\label{sec:intro-constraints}

With linear constraints the optimizer lies on an unknown face, and finding
that face is a search problem separate from the arithmetic.
Section~\ref{sec:constraints} treats globally strongly convex quartics on
rational polyhedra and, where stated, strongly monotone cubic variational
inequalities. For an arbitrary rational polyhedron, exact comparison of the
value, a coordinate, or an explicit observable lies in
$\mathrm{UP}^{\PosSLP}\cap\mathrm{coUP}^{\PosSLP}$: the full active set is
a unique certificate, and the normal-cone condition is verified by a linear
program whose objective is given by a circuit. Whether the comparison lies in
$\mathrm P^{\PosSLP}$, equivalently whether one $\PosSLP$ instance suffices, is
open. The section removes the search in four ways.
\begin{itemize}
\item For globally strongly convex objectives of unary degree, a supplied
  lower bound on all nonzero optimal slacks lets ordinary approximation
  recover the active set, after which one $\PosSLP$ instance decides each
  relation (Proposition~\ref{prop:upper-slack-gap} in
  Section~\ref{sec:upper}).
\item If the cubic and quartic parts of $f$ depend on only $k$ linear forms,
  exact comparison and recovery of the full active set take ordinary
  deterministic time $F(k)L^C$, with $F(k)=2^{O(k)}$ and an absolute constant
  $C$. The quadratic part, the number of constraints, and their rank are
  unrestricted. The optimizer is returned as an affine image of the unique
  stationary point of a strongly convex quartic in at most $k$ variables.
\item For objectives of fixed degree on boxes whose Hessian interactions form
  a forest, have nonpositive off-diagonal entries, or have a positive
  definite comparison matrix, and on network-flow polyhedra with separable
  objectives, exact constrained Newton steps can be executed with
  polynomially many arithmetic operations. Combined with the compilation
  theorem, each exact predicate about the optimizer, and each polynomial-size
  Boolean combination of such predicates, reduces to one $\PosSLP$ instance
  (Corollary~\ref{cor:structured-single-sign}). The full active set is
  recovered by a polynomial list of such instances, prepared without
  adaptivity.
\item If the constraint normals have rank $r$, a Las Vegas algorithm with a
  $\PosSLP$ oracle runs in expected time $(r+1)^{O(r)}L^C$.
\end{itemize}

Integer variables add a second search. Let the quartic objective be strongly
convex in all variables jointly, and let $t$ variables be integer. An
ordinary algorithm running in time $2^{O(t\log(t+1))}L^C$ lists candidate
integer blocks that include every optimal block when the continuous fibers
are unconstrained; there are at most $2^t$ optimal blocks. Exact selection
among the candidates then uses $\PosSLP$ queries that can be prepared in
advance. Under arbitrary mixed linear constraints the list takes ordinary
time $a(t)L^C$ for a computable function $a$, and the methods above combine
to give an ordinary algorithm in time $F(k+t)L^C$, or a Las Vegas algorithm
in expected time $F(t,r)L^C$ with a $\PosSLP$ oracle. Thus discrete search
needs only ordinary precision, and exact arithmetic is needed only to select
among the candidates. This selection cannot be avoided: deciding the value of
a single binary variable at the unique optimum remains $\PosSLP$-hard even
when the optimal value is known, the optimizer is rational and bounded, and
the continuous constraints have rank one.

\subsection{Quadratic and algebraic arithmetic interfaces}
\label{sec:intro-further-arithmetic}

Quadratic constraints change the arithmetic even when they are convex.
Appendix~\ref{app:qc} studies positive semidefinite quadratic constraints
whose Hessians span a space of dimension $h$ (denoted $k$ in that appendix).
The optimal value and a
canonical optimizer lie in one number field of degree at most
$\max_{0\le s\le\min(h,n)}2^s\binom ns$, and the degree bound is sharp.
For fixed $h$, exact feasibility and common-field point recovery take
ordinary polynomial time, but expanded minimal-polynomial output admits no
bound $F(h)L^C$ with an absolute exponent $C$. The rational-witness boundary
is also sharp: nonempty systems of Hessian span at most two have rational
points of polynomial bit length, whereas three ellipsoids can intersect at a
single irrational point. The appendix also treats one-field optimizer
certificates, rational positive-aggregate infeasibility certificates, and
linear and convex quadratic programs with rational matrices and affine data
in an explicitly represented real number field. These results explain which
parts of the higher-degree account already arise from quadratic constraints.

Appendix~\ref{app:further-arithmetic} extends the description and input
questions in two directions. For rational quadratic constraints that need
not be convex, with Hessian span $h$, a feasible point has a common-field
algebraic description of degree and total length $L^{O(h+1)}$. This
feasibility bound follows from the quadratic-map sampling theorem of
Grigoriev and Pasechnik~\cite[Theorem~1.2]{GrigorievPasechnik2005} and an
elementary argument on the minimal face of the affine constraints. The
appendix controls boxed optimal values, finite infima that need not be
attained, and minimum-norm optimizers when attainment holds. It keeps affine
constraints in the optimization formulas: an affine inequality that is
inactive at an optimizer can still change the global optimal value when
deleted.
For fixed $h$, feasibility belongs to $\NP$, and the status, finite value,
and an attained optimizer can be recovered in $\mathrm{FP}^{\NP}$.
These are exact arithmetic guarantees, rather than ordinary polynomial-time
algorithms for nonconvex optimization.

The second extension allows cone data in one explicitly represented number
field $K\subseteq\R$, with its real embedding specified. With fixed
integer dimension $t$ and a fixed dimension $h$ for the span over $K$ of the
continuous Hessians of the squared cone residuals, exact feasibility and
witness recovery take ordinary polynomial time,
without a supplied box or a Slater condition. The continuous witness is
written in the joint field generated by $K$ and its coordinates. For a
continuous system its minimum-norm point has absolute field degree and total
description length $L^{O(h+1)}$; in the mixed-integer case the algorithm first
finds an integer assignment and then applies this result to its exact
continuous fiber. Known algebraic
thresholds can therefore be used exactly in a later cone problem. The
extension concerns explicit algebraic input, sign and rounding operations,
and common-field output; its sampling and fixed-dimensional integer
ingredients are prior methods.
Exact second-order cone feasibility with unrestricted $h$ is
$\PosSLP$-hard (Remark~\ref{rem:models-socp}); in those gate constructions
$h$ can grow with the circuit size. The polynomial-time bound here fixes
both $t$ and $h$.

\subsection{Exact optimizers}
\label{sec:intro-optimizers}

Exact optimizers can be irrational even in very small, well-conditioned
instances. With
\[
A=12599x^2-10000xy+7937y^2-15874x-12599y+20000,
\]
the integer quartic
\[
F(x,y)=A^2+10000\bigl[(x^2-y)^2+(y^2-2x)^2\bigr]
\]
satisfies $\nabla^2F\succeq4124\,I$ on $\R^2$, has a positive definite
integer full Hessian Gram, and has the single zero $(2^{1/3},2^{2/3})$
(Proposition~\ref{prop:algebraic-bivariate}). Hence the constraint $F\le0$ is
satisfiable but has no rational solution. This answers negatively the
rational-witness question for convex quartics that is marked open in the
arXiv version of~\cite{SlotSteurerWiedmer2025}. It does not settle the
complexity of the decision, since the point has a short algebraic
description; Section~\ref{sec:intro-exact} shows that the decision is
$\PosSLP$-hard. A univariate convex quartic with rational minimum has a
rational minimizer~\cite[Lemma~C.3]{SlotSteurerWiedmer2025}, so two
variables are needed.

Theorem~\ref{thm:singleton-field} identifies which numbers occur. A real
algebraic number is a coordinate of a point isolated by finitely many
rational convex polynomial inequalities if and only if its minimal polynomial
has exactly one real root; the same numbers are the coordinates of unique
minimizers of rational convex polynomials on $\R^n$, and of single real zeros
of rational polynomials. Conversely, each such irrational number is a
coordinate of the unique zero of a single rational, strongly convex quartic
that is a sum of squares and carries a positive definite rational full
Hessian Gram; the quartic is constructed in polynomial time from the dense
minimal polynomial. For the prescribed point of consecutive powers of the
number, $(d+1)/2$ variables suffice and are necessary, where $d$ is the
degree. The restriction to isolated points matters: minimizing $-x$ subject
to $x^2\le2$ has the unique optimizer $\sqrt2$, whose minimal polynomial has
two real roots.

Certified strong convexity does not bound the algebraic degree.
Section~\ref{sec:algebraic} constructs integer quartics in $n$ variables with
$O(n^2)$ monomials, coefficients of $O(\log n)$ bits, $n+1$ integer square
factors, and a positive definite integer full Hessian Gram, whose unique zero
$p$ satisfies
\[
[\Q(p):\Q]=[\Q(p_1):\Q]=d_n=\frac{2^{n+1}-(-1)^{n+1}}{3}.
\]
The Hessian condition number at the zero is $O(n^2)$. After a rational
translation, the minimal polynomial of the first coordinate has all $d_n+1$
coefficients nonzero. Even with the full Hessian Gram counted, the input
length is $L=O(n^4\log n)$, so dense and sparse coefficient lists of this
minimal polynomial have length $2^{\Omega(n)}=2^{\Omega((L/\log L)^{1/4})}$:
exponential in the dimension and superpolynomial, but not exponential, in
$L$. Every coordinate of the zero is nevertheless an integer power of
$2^{1/d_n}$ whose exponent has $O(n)$ bits, so one real root gate, whose
index $d_n$ is written in binary, followed by $O(n)$ multiplications and
divisions computes the whole optimizer. In the other direction, if a rational sum of
squares of quadratics has a unique real zero with positive definite Hessian
there, the degree of the field generated by the coordinates of the zero is
odd and at most $2^n-1$. Global convexity improves the bound to $2^n-3$ for
$n\ge3$ and $2^n-5$ for $n\ge4$, and the exact maxima for $n=2,3,4,5$ are
$3,5,11,21$, attained by the construction above.
Theorem~\ref{thm:sos-length} proves that a globally convex polynomial of
degree exactly four that is a sum of squares and has a zero with positive
definite Hessian needs at least $n+1$ squares, even with real coefficients;
the bound is sharp and our constructions attain it. Auxiliary variables also
matter: some algebraic numbers whose minimal polynomials have $O(b)$ bits
require degree $2^{\Omega(b)}$ in every rational convex univariate
realization, while two-variable quartics of polynomial size realize them.

Exact rational optimizers and rational feasible points can be long.
Theorem~\ref{thm:rational-height} gives strongly convex quartics in
$2(k+1)$ variables, with short rational square factors and a full positive
definite Hessian Gram, whose rational minimizer lies in a fixed box and has
last coordinates with reduced denominator exactly $5^{2^k}$. A rescaled
variant has a minimizer whose last denominators are divisible by $5^{2^k}$
and a Hessian condition number at the minimizer that is polynomial in $k$;
the two properties are proved for these two related families, not for one
instance. Section~\ref{sec:heights} also gives certified quartics
in $n$ variables whose zero sublevel set is compact with nonempty interior,
yet every rational point of it needs $\Omega(n2^{n/2})$ denominator bits.
These bounds are exponential in the dimension, hence superpolynomial in the
input length, which is polynomial in $n$. They concern expanded rational
output. The long rational minimizers above are values of polynomial-size
circuits built by repeated squaring, and when the minimum of an instance of
Theorem~\ref{thm:exact-upper} is negative, the Newton circuit of that theorem
gives a strictly feasible point as a polynomial-size shared rational circuit
(Section~\ref{sec:upper}). Expanded formats can also be long for
well-conditioned instances. For quartics that are strongly convex on a box,
with bounded curvature and a path as interaction graph,
Section~\ref{sec:points} shows that the minimal polynomial of an optimizer
coordinate can have a number of nonzero coefficients exponential in the
dimension, and that every rational box certifying the optimizer inside the
feasible box can need endpoints whose bit length is exponential in the
dimension. With interaction graphs of treewidth two, the same holds for
rational boxes on which the sign of an active multiplier is certified. In each
of these families a short recurrence describes the optimizer exactly.

\subsection{Exact certificates}
\label{sec:intro-certificates}

Certificates of optimality have their own arithmetic. Let $f$ be a quartic
with a positive definite full Hessian Gram and positive minimum. A positive
definite rational polynomial Gram matrix of total bit length
$\poly(L)2^{O(n)}$ exists, and Section~\ref{sec:heights} shows that
$\Omega(n2^{n/2})$ bits can be necessary, even though the same instances have
short singular rational Gram matrices. Theorem~\ref{thm:circuit-gram}
constructs, in deterministic polynomial time and without any sign oracle, a
polynomial-size shared rational circuit for a Gram matrix of $f$ that is
positive definite exactly when the minimum is positive. Deciding whether this
matrix is positive definite is therefore equivalent to deciding whether the
minimum is positive, which is $\PosSLP$-complete in this class. With a
positive definite full Hessian Gram, the optimal order-two moment matrix is
unique, and the field $\Q(p)$ generated by the coordinates of the minimizer is
the least field over which a maximal-rank optimal Gram matrix, or a nonzero
matrix exposing the face of optimal Gram matrices, can be written. For the
rational minimizers of large height above, the optimal moment matrix, every
maximal-rank rational optimal Gram matrix, and every nonzero rational exposing
matrix have entries whose bit length is exponential in the dimension, while
short optimal Gram matrices of lower rank remain.

Rational certificates at the optimum can fail to exist.
Section~\ref{sec:fields} gives a ternary integer quartic $F_\star$
(Example~\ref{ex:fields-ternary}) with $\nabla^2F_\star\succeq I$, a positive
definite rational full Hessian Gram, and minimum zero, for which a
sum-of-squares decomposition, or a positive semidefinite Gram matrix, with
entries in a real field $E$ exists exactly when $2^{1/5}\in E$. Three variables
are the minimum under this certificate hypothesis, and $F_\star+\eta$ is a
rational sum of squares for every rational $\eta>0$. For every prime $\ell\ge5$
there is an analogous integer quartic in $(\ell+1)/2$ variables, with
coefficients and Hessian Gram entries of $O(\log\ell)$ bits, whose least
field is $\Q(2^{1/\ell})$; it is constructed in time polynomial in $\ell$. A
family of quartics in $3k$ variables, of size polynomial in $k$, has least
field $\Q(2^{1/5^k})$; in every algebraic certificate some individual
coefficient has degree at least $5^k$. For this family, compact formats avoid
the cost: the certificate coefficients are outputs of one polynomial-size
circuit with $k$ real fifth-root gates, and adjoining the coordinates of the
minimizer as auxiliary variables gives a polynomial-size rational certificate
modulo rational quadratic equations that these coordinates satisfy.
Section~\ref{sec:fields} also proves a sufficient condition for rational
descent. Suppose that some rational globally convex quartic that is a rational
sum of squares vanishes at a point $p$ with positive definite Hessian there,
that the rational polynomials of degree at most two vanishing at $p$ span a
space of dimension $n+1$, and that their pairwise products span the rational
polynomials of degree at most four that vanish at $p$ together with their
gradients. Then every rational globally convex quartic with minimum zero at
$p$ and positive definite Hessian there is a rational sum of squares.
For the cyclic points of Section~\ref{sec:algebraic}, the dimension and
product-independence hypotheses hold for $n\ge4$. Whether the products
exhaust the stationary quartics for all these points remains open.

The format of a certificate changes what exists and how large it is.
Section~\ref{sec:fields-certificates} shows that the explicit ternary quartic
$F_\star$ has a short exact rational certificate as a sum of squares of
rational functions with the common denominator $1+\norm{x}^2$, and that the
same holds for the tower quartics, whose polynomial certificates need a
coefficient of degree $5^k$, after a polynomial-size increase of a scale
parameter; two is the least possible denominator degree. For a scaled family,
every fixed power $(1+\norm{x}^2)^N$ fails as a multiplier for large scale.
The least admissible power, which is finite, grows as
$\Omega(\log L/\log\log L)$, while a multiplier adapted to the scale keeps the
certificate of size $O(L)$. For a sum $f(x)+g(y,z)$ of strongly convex
quartics in disjoint variables with zero minima, a short joint rational sum of
squares can coexist with the absence of any rational certificate separated by
blocks, and the separated least field can be $\Q(2^{1/5^k})$. After a
certified quartic with a tiny positive minimum is added to the second block,
separated rational certificates exist but need $\Omega(k2^k)$ bits.
Theorem~\ref{thm:quadratic-graph} supplies the tool that makes the second
block strongly convex: from a quartic with zero minimum and a positive
definite full Hessian Gram, it constructs a polynomial-size certified strongly
convex quartic whose unique zero is the optimizer together with all products
of its coordinates.

\subsection{Summary of results}
\label{sec:intro-summary}

Table~\ref{tab:intro-results} lists the main results by output. In the model
column, ``ordinary'' means a deterministic algorithm in the bit model without
an oracle. ``One instance'' means a polynomial-time many-one reduction to
$\PosSLP$, and ``complete'' means that, in addition, $\PosSLP$ reduces to the
problem. ``Conditional'' marks a reduction from Square Root Sum or $\PosSLP$
to an approximation task.

{\small
\setlength{\tabcolsep}{4pt}
\setlength{\LTcapwidth}{\textwidth}
\begin{longtable}{@{}>{\raggedright\arraybackslash}p{0.17\textwidth}
  >{\raggedright\arraybackslash}p{0.27\textwidth}
  >{\raggedright\arraybackslash}p{0.35\textwidth}
  >{\raggedright\arraybackslash}p{0.15\textwidth}@{}}
\caption{Main results by output. Promised structure is part of the setting;
$L$ is the input length, $D$ the degree, $q$ the requested precision,
$\Lambda$ an upper bound on the coordinate second derivatives in the core, and
$\sigma$ the tilt width.}
\label{tab:intro-results}\\
\toprule
Output & Setting & Result and model & Where\\
\midrule
\endfirsthead
\toprule
Output & Setting & Result and model & Where\\
\midrule
\endhead
\bottomrule
\endfoot
\multicolumn{4}{@{}l}{\emph{Values and points}}\\*
Fixed optimizer & Globally convex, explicit sparse, degree $D$; rational
polyhedron & Error exponent $1/D$, constant of $\poly(L,D)$ bits; time
$\poly(L,D,q)$; ordinary & Thm.~\ref{thm:global-point}\\
Fixed optimizer & Cubic, convex on a bounded rational polytope & Exponent
$1/4$, constant of $\poly(L)$ bits; time $\poly(L+q)$; ordinary &
Thm.~\ref{thm:cubic-point}\\
Some optimizer, accuracy $1/4$ & Quartic convex on a box, unique optimizer &
Decides Square Root Sum (treewidth at most two) or $\PosSLP$; conditional &
Thm.~\ref{thm:points-quartic-lower}\\
Active bound & Strongly convex cubic on a box & Decides Square Root Sum, while
points are easy; conditional & Prop.~\ref{prop:points-active}\\
Full selected optimizer & Residual-convex cubic on a product box, one random
core tilt & Expected $k^{O(k)}(1+\Lambda/\sigma)^k\poly(L+q)$, correct for
every draw & Thm.~\ref{thm:recourse-cubic}\\
Values and fixed core & Fixed degree; residual convexity on a product box,
or a supplied convexifier on a coupled polytope & One finite core-noise law;
every draw correct; one expected-work factor for all precisions &
Thm.~\ref{thm:recourse-core}\\
Exact optimizer and value & Quadratic in these core-noise models & Rational
output of polynomial length; expected parameterized work &
Cor.~\ref{cor:recourse-qp}\\
Full selected optimizer & Supplied convexifier; cubic on a polytope or
globally convex after convexification & Expected parameterized work; ordinary
polynomial expected work when jointly convex &
Thm.~\ref{thm:recourse-convexified}, Cor.~\ref{cor:recourse-joint}\\
\multicolumn{4}{@{}l}{\emph{Exact decisions}}\\*
Sign or equality of $h(p)$ & Explicit $f$, $h$, unary degree,
$\nabla^2f\succeq\mu I$ supplied & One instance per relation, no gap
assumed & Thm.~\ref{thm:exact-upper}\\
Same, at a zero of a map & Strongly monotone polynomial map, unary degree &
One instance per relation; order tests complete for certified cubic maps,
hard already for gradient maps; equality only upper &
Thm.~\ref{thm:upper-monotone}\\
Order tests of value and coordinates; nonnegativity; real SOS & Quartic with
checked full Hessian Gram $A\succ0$ & Complete; rational SOS membership hard;
equality only upper & Thm.~\ref{thm:quartic-complete}\\
Coordinate order & As above, rational minimizer in $[-1,1]^n$, minimum known
to be $0$ & Complete & Thm.~\ref{thm:rational-optimizer}\\
Feasibility of $f\le0$ & Quartic with checked $A\succ0$, on a halfspace or
the unit cube & $\PosSLP$-hard; in
$\mathrm{UP}^{\PosSLP}\cap\mathrm{coUP}^{\PosSLP}$ &
Cor.~\ref{cor:reductions-feasibility},
Thm.~\ref{thm:constraints-unambiguous}\\
Value, coordinate, observable & Strongly convex quartic on a rational
polyhedron & $\mathrm{UP}^{\PosSLP}\cap\mathrm{coUP}^{\PosSLP}$ &
Thm.~\ref{thm:constraints-unambiguous}\\
Same & Strongly convex polynomial on a rational polyhedron, supplied gap for
nonzero slacks & One instance per relation; active set by ordinary
approximation & Prop.~\ref{prop:upper-slack-gap}\\
Same, and active set & Nonlinear dimension $k$ & Ordinary $F(k)L^C$ &
Thm.~\ref{thm:constraints-nonlinear}\\
Any Boolean predicate & Fixed degree; forest, nonpositive off-diagonal, or
comparison-matrix Hessian on a box; separable network flow & One instance &
Cor.~\ref{cor:structured-single-sign}\\
Value, coordinate, observable & Constraint rank $r$ & Las Vegas
$(r+1)^{O(r)}L^C$ with $\PosSLP$ oracle & Thm.~\ref{thm:constraints-rank}\\
Optimal integer blocks & $t$ integer variables & Ordinary candidate lists in
FPT time; exact selection by $\PosSLP$ queries &
Thms.~\ref{thm:constraints-candidates}, \ref{thm:constraints-mixed-candidates}\\
One binary decision & Known optimal value, rational optimizer, rank one &
$\PosSLP$-hard & Cor.~\ref{cor:constraints-binary}\\
\multicolumn{4}{@{}l}{\emph{Exact optimizers}}\\*
Possible coordinates & Points isolated by rational convex polynomial
inequalities & Exactly the numbers with one real conjugate; one certified
quartic realizes each, in polynomial time & Thm.~\ref{thm:singleton-field}\\
Coordinate field degree & Certified convex rational SOS quartics, unique
nondegenerate zero & Degree $d_n=(2^{n+1}-(-1)^{n+1})/3$ with $O(\log n)$-bit
integer data; at most $2^n-5$ for $n\ge4$; exact maxima for $n\le5$ &
Thms.~\ref{thm:algebraic-cyclic}, \ref{thm:algebraic-degree-bounds}\\
Number of squares & Convex quartic SOS, nondegenerate zero & At least $n+1$
real squares; sharp & Thm.~\ref{thm:sos-length}\\
Rational minimizer & Certified quartics in $n$ variables & Last coordinates
with denominator exactly $5^{2^{n/2-1}}$ & Thm.~\ref{thm:rational-height}\\
Rational feasible points & Certified quartics in $n$ variables, sublevel set
with interior & Every rational point needs $\Omega(n2^{n/2})$ bits &
Thm.~\ref{thm:heights-witness}\\
\multicolumn{4}{@{}l}{\emph{Exact certificates}}\\*
Interior Gram matrix & Positive minimum, full Hessian Gram & $\poly(L)2^{O(n)}$
bits suffice, $\Omega(n2^{n/2})$ can be needed; shared rational circuit in
polynomial time & Thms.~\ref{thm:interior-gram}, \ref{thm:circuit-gram}\\
Moment and exposing matrices & Full Hessian Gram & Unique moment optimum;
least field $\Q(p)$; long entries for long rational minimizers &
Thm.~\ref{thm:heights-moment}, Cor.~\ref{cor:heights-gram}\\
SOS coefficient field & Certified quartics with minimum $0$ & Least field
$\Q(2^{1/\ell})$ in $(\ell+1)/2$ variables; $\Q(2^{1/5^k})$ in $3k$ variables,
with a coefficient of degree $\ge5^k$; polynomial-size root-circuit formats
for this family & Thms.~\ref{thm:fields-prime}, \ref{thm:fields-tower},
\ref{thm:fields-individual}, \ref{thm:fields-compact}\\
Certificate format & Rational functions, radial multipliers, block
separation & Quadratic common denominators suffice, also for the tower after
a scale increase; growing radial order; separation fails or costs
$\Omega(k2^k)$ bits & Cor.~\ref{cor:fields-tower-denominator},
Thms.~\ref{thm:fields-radial}, \ref{thm:fields-radial-height},
\ref{thm:fields-blocks}, \ref{thm:fields-blocks-height}\\
\multicolumn{4}{@{}l}{\emph{Further arithmetic interfaces}}\\*
Feasible point, finite infimum, attained optimizer & Rational quadratics,
not necessarily convex; constraint Hessian span $h$ & Common-field algebraic
descriptions of degree and total length $L^{O(h+1)}$; finite infima may be
unattained; fixed-$h$ feasibility in $\NP$ and exact output in
$\mathrm{FP}^{\NP}$ & Thms.~\ref{thm:ncq-feasible},
\ref{thm:ncq-infimum}, \ref{thm:ncq-oracle}\\
Cone feasibility and witnesses & Explicit common real number field; integer
dimension $t$; continuous Hessians of squared cone residuals span $h$
directions over that field &
Ordinary polynomial time for fixed $t,h$; continuous point in a joint field
with degree and total length $L^{O(h+1)}$ for continuous input; no box or
Slater hypothesis &
Thms.~\ref{thm:cone-feasibility}, \ref{thm:cone-witness}\\
\end{longtable}
}

\subsection{Relation to prior work}
\label{sec:intro-prior}

\paragraph{Convex polynomial optimization and error bounds.}
Slot, Steurer, and Wiedmer establish, for globally convex polynomials with
exponents written in unary on rational polyhedra, the dichotomy between
unboundedness and attainment together with a bound on the norm of an
optimizer whose logarithm is polynomial in the input
length~\cite[Theorem~1.1]{SlotSteurerWiedmer2025}, and polynomial-time value
approximation~\cite[Corollary~1.2]{SlotSteurerWiedmer2025}.
Section~\ref{sec:points} uses the same structural facts, which rest on their
analysis and on a tangent quadratic bound from the proof of Lemma~5 of
Ahmadi, Chaudhry, and Zhang~\cite{AhmadiChaudhryZhang2024}, and rederives
attainment and the radius with explicit rational constants. Global error
bounds for convex polynomials, also relative to polyhedral constraints, are
studied in~\cite{Li2010,Li2013,Yang2009,Ngai2015}; the bounds
of~\cite{Li2010,Li2013,Ngai2015} are H\"older-type estimates whose constants
depend on the function. The addition of Theorem~\ref{thm:global-point} is
the quantitative selector: an error-bound constant computed from the input,
of polynomial bit length on arbitrary rational polyhedra, evaluation of
sparse polynomials without expanding affine substitutions, and a
regularization schedule that approximates one fixed optimizer at every
precision. The form of the error exponent is not claimed as new. The
distinction between approximating the residual of a solution and
approximating the solution itself, and the hardness of the latter for
equilibria, are due to Etessami and Yannakakis~\cite{EY2010};
Section~\ref{sec:points} shows the same phenomenon for minimizers of
polynomials that are convex on a box. Recognizing convexity is NP-hard for
quartics on $\R^n$~\cite{AhmadiOlshevskyParriloTsitsiklis2013} and for cubics
on a box~\cite{AhmadiHall2020}. Convexity is therefore a promise in our
theorems, unless a certificate is supplied and checked.

\paragraph{Exact arithmetic, Newton circuits, and $\PosSLP$.}
Allender, B\"urgisser, Kjeldgaard-Pedersen, and
Miltersen~\cite{AllenderEtAl2009} introduce $\PosSLP$ and show that
polynomial-time constant-free real computation on Boolean inputs corresponds
to $\mathrm P^{\PosSLP}$~\cite[Proposition~1.1]{AllenderEtAl2009}; divisions
are removed there by computing numerators and denominators separately, and
Newton iteration gives short circuits that approximate square
roots~\cite[Section~1.4]{AllenderEtAl2009}. Etessami, Stewart, and
Yannakakis~\cite[Appendix~C]{EtessamiStewartYannakakis2012} reduce threshold
comparisons at least fixed points of probabilistic polynomial systems to one
$\PosSLP$ instance, by storing polynomially many Newton steps in a circuit
and combining them with an algebraic separation bound. Our upper reduction
uses this architecture for a different class of inputs, globally strongly
convex explicit polynomials and strongly monotone polynomial maps. The
class-specific steps are a polynomial-bit starting point obtained from
convex value approximation, a separation bound that holds even when the
complex critical set of the objective has positive dimension, the treatment
of equality, and the handling of supplied certificates. For the compilation
theorem the nearest prior work is the oracle characterization of
$\mathrm P^{\PosSLP}$ through constant-free real
computation~\cite{AllenderEtAl2009,BuergisserJindal2024} and specific
many-one reductions such as that of~\cite{EtessamiStewartYannakakis2012}.
These works state adaptive (Turing) relationships or particular reductions,
and the sign iteration and division elimination are classical numerical
tools. Section~\ref{sec:upper} gives the explicit $O(S^2)$ compilation with
its error budget and interpreter in full; we make no priority claim for it.
B\"urgisser and Jindal~\cite{BuergisserJindal2024} study the hardness of
$\PosSLP$ itself; no unconditional NP-hardness of $\PosSLP$ is known, so our
$\PosSLP$-hardness results are not NP-hardness results.

\paragraph{Exact conic optimization.}
Tarasov and Vyalyi~\cite[Theorems~3 and~4]{TarasovVyalyi2008} reduce
$\PosSLP$ to exact semidefinite feasibility. They compare the values of two
circuits built from the constant $1$ by the operations $x+y$ and $x^2/2$,
represent each value as the optimum of a semidefinite program, and use
duality. Gate constraints on two-by-two matrices are rotated second-order
cone constraints, and standard conic duality then gives the same hardness for
exact second-order cone feasibility (Remark~\ref{rem:models-socp}). Hardness
of exact convex optimization in general is therefore not new. Our lower
bounds hold in a much smaller class: a single explicit quartic objective,
globally strongly convex with a checkable strict certificate and without
constraints, for which Theorem~\ref{thm:exact-upper} supplies the matching
upper bound. Squaring the residuals of the gate equations does not preserve
convexity, so this restriction requires the realization of
Lemma~\ref{lem:quartic-realization}.

\paragraph{Algebraic degree and irrational optimizers.}
Nie and Ranestad~\cite{NieRanestad2009} compute the algebraic degree of
generic polynomial optimization problems. Such counts concern generic data,
and an optimizer of a particular instance can have much smaller degree.
Irrational singletons defined by convex polynomials were known before. In the
arXiv version of~\cite{SlotSteurerWiedmer2025}, Appendix~C gives a univariate
convex sextic whose minimum zero is attained at a single irrational point.
Bienstock, Del Pia, and
Hildebrand~\cite[Example~2.1]{BienstockDelPiaHildebrand2023} give a cubic
whose nonpositive sublevel set within a rational rectangle is the single point
$(2^{1/3},2^{2/3})$; its Hessian is positive definite on that rectangle, so
the cubic is convex there, but no cubic is convex on all of $\R^2$. Our
degree results concern specific certified convex classes. They characterize
which numbers occur, construct a certified strongly convex quartic for each
of them from its dense minimal polynomial, bound the degree from above,
determine the maxima in dimensions at most five, and exhibit degree
exponential in the dimension with coefficients of logarithmic bit length.

\paragraph{Rational sums of squares.}
Scheiderer~\cite{Scheiderer2016} constructs rational forms that are sums of
squares over $\R$ but not over $\Q$, shows that for every number field with a
real embedding some rational form is a sum of squares over $\R$ but not over
that field~\cite[Corollary~2.11]{Scheiderer2016}, and characterizes the
nonnegative rational ternary quartic forms that are not sums of squares over
$\Q$~\cite[Theorem~4.1]{Scheiderer2016}. Sums of squares descend from totally
real number fields to $\Q$~\cite[Theorem~1.4]{Hillar2009},
\cite[Theorem~1.2]{Scheiderer2016}; the least fields $\Q(2^{1/\ell})$ and
$\Q(2^{1/5^k})$ in our examples are not totally real, so this descent does not
apply to them. Laplagne~\cite[Section~3]{Laplagne2024} gives rational quartic
forms that are sums of squares over $\Q(2^{1/3})$ but not over $\Q$, one of
them positive away from the origin; convexity plays no role in those
constructions. Chua, Plaumann, Sinn, and
Vinzant~\cite{ChuaPlaumannSinnVinzant2017} record that rational positive
semidefinite Gram matrices and rational sums of squares are equivalent, and
that the analogous statement fails over general ordered fields. Peyrl and
Parrilo~\cite{PeyrlParrilo2008} round interior numerical Gram matrices to
exact rational certificates. A sum-of-squares-convex polynomial that is
stationary at a zero is a real sum of squares, by integrating its Hessian
certificate~\cite{HeltonNie2010,AhmadiParrilo2013,Lasserre2009}. Our field
results add global strong convexity with a strict rational convexity
certificate, attained minimum zero, exact least fields for both sums of
squares and Gram matrices, field degree $5^k$ for inputs of size polynomial in
$k$ together with a lower bound for individual coefficients and
polynomial-size root-circuit certificates for that family, and a sufficient
condition for rational descent.

\paragraph{Size of certificates.}
Algorithms that produce exact rational sum-of-squares certificates have bit
bounds that depend on a supplied eigenvalue margin or on the existence of a
rational certificate~\cite{PeyrlParrilo2008,SafeyElDinZhi2010,
MagronSafeyElDin2021,DavisPapp2022,GaertnerMagronVallentin2026}, and bit
complexity lower bounds are known for sum-of-squares proofs from
constraints~\cite{RaghavendraWeitz2017}. Our size bounds concern unconstrained
certificates for a single short certified quartic, and each is stated for a
specific format: interior Gram matrices, maximal-rank optimal Gram matrices,
exposing matrices, prescribed multipliers, or certificates separated by
blocks. They identify certified instances where interior certificates are
necessarily long and where rational certificates at the optimum do not
exist. Kojima, Kim, and Waki~\cite{KojimaKimWaki2005} discuss real sums of
squares respecting a block separation; our block results concern rational
coefficients and their size.

\paragraph{Constraints and integer variables.}
The ingredients of Section~\ref{sec:constraints} include essential-variable
extraction~\cite{Carlini2006}, approximation algorithms for objectives with a
low-rank nonlinear part~\cite{MittalSchulz2013,KannanRademacher2009},
quadratic convergence of exact proximal Newton
steps~\cite{LeeSunSaunders2014}, exact algorithms for structured
box-constrained and network-flow quadratic
programs~\cite{PangHan2023,Vegh2016}, violator
spaces~\cite{GaertnerMatousekRuestSkovron2008}, convex integer minimization
in fixed dimension~\cite{HildebrandKoeppe2013}, integer feasibility
algorithms that query separation oracles only at integer
points~\cite{AriHildebrand2026,HildebrandGoess2024,Basu2023Integers}, and
lattice basis reduction~\cite{LenstraLenstraLovasz1982}. Our additions there
are the exact arithmetic interfaces that these methods need for algebraic
optimizers and their compositions. For example, the integer-query theorems
are not applied directly to the algebraic value of a continuous fiber;
Section~\ref{sec:constraints} replaces them by its own approximate
projected-gradient and integer-only cutting argument.

\paragraph{Quadratic sampling and algebraic cone input.}
The feasibility bound in Appendix~\ref{app:further-arithmetic} is a
polyhedral consequence of Grigoriev and Pasechnik's
component-sampling theorem~\cite[Theorem~1.2]{GrigorievPasechnik2005}.
Its sampling architecture is prior work. The appendix gives separate
proofs for finite infima and attained minimum-norm optimizers, using ordered
regularization and perturbation limits and controlling one joint field for
the selected limiting tuple. Its cone algorithm uses that arithmetic
interface together with the integer-witness method of Khachiyan and
Porkolab~\cite[Theorem~1.1]{KhachiyanPorkolab2000} and rational mixed-integer linear
feasibility at fixed integer dimension~\cite[Section~5]{Lenstra1983}.
The extension specifies one real embedding for explicit algebraic input,
preserves exact feasibility through rational rounding and cone lifts, and
returns continuous coordinates in a field containing the input field.
These input, formula, sign, rounding, and output guarantees are distinct
from the established sampling and integer algorithms on which they rely.

\subsection{Scope of the claims}
\label{sec:intro-scope}

The hardness results for continuous convex polynomial families are reductions
from $\PosSLP$ or Square Root Sum. They do not establish NP-hardness or
unconditional running-time lower bounds. Appendix~\ref{app:further-arithmetic}
also proves strong NP-hardness of nonconvex quadratic feasibility and
attainment at a supplied zero infimum, even with one constraint-Hessian
direction. Upper bounds relative to $\PosSLP$ are not
ordinary polynomial-time algorithms, since the constructed circuits compute
integers with exponentially many bits. $\PosSLP$ does not describe polynomial
optimization in general: our upper bounds require explicit input of unary
degree and a global strong convexity bound, or one of the constrained
structures of Section~\ref{sec:constraints}. Three questions about exact
decisions remain open: whether exact comparison for globally strongly convex
quartics over arbitrary rational polyhedra lies in $\mathrm P^{\PosSLP}$,
equivalently reduces to one $\PosSLP$ instance, rather than only in
$\mathrm{UP}^{\PosSLP}\cap\mathrm{coUP}^{\PosSLP}$; whether exact comparison
at minimizers of convex quartics that are not strongly convex has any upper
bound relative to $\PosSLP$; and whether equality tests are $\PosSLP$-hard.
Lower bounds on output size are stated for specific expanded formats and do
not exclude compact descriptions. Most of them concern families whose input
length is polynomial in the dimension and are exponential in the dimension,
hence superpolynomial in the input length; none of these is claimed to be
exponential in the input length. Promised properties, such as convexity,
strong convexity, a unique real root, a rational optimizer, a zero minimum,
or a slack gap, are not recognized by our algorithms unless a certificate is
supplied and checked.

\subsection{Organization}
\label{sec:intro-organization}

Section~\ref{sec:models} fixes encodings, promises and certificates, output
contracts, and computational models, and proves a few elementary facts used
throughout. Section~\ref{sec:points} studies value-to-point conversion.
Sections~\ref{sec:upper} and~\ref{sec:reductions} prove the upper and lower
bounds for exact decisions, and Section~\ref{sec:constraints} extends them to
constraints and integer variables. Sections~\ref{sec:algebraic}
and~\ref{sec:heights} study the algebraic degree and the bit length of exact
optimizers and of interior, moment, and exposing matrices.
Sections~\ref{sec:fields} and~\ref{sec:fields-certificates} study
coefficient fields and formats of sum-of-squares certificates.
Section~\ref{sec:recourse} treats cubic recourse with a nonconvex core.
Section~\ref{sec:discussion} discusses consequences for exact optimization
and open questions.

The main sections state results, explain why they matter, and outline the
proofs. The complete proofs for
Sections~\ref{sec:points}--\ref{sec:recourse} are in
Appendices~\ref{app:points}, \ref{app:upper}, \ref{app:reductions},
\ref{app:constraints}, \ref{app:algebraic}, \ref{app:heights},
\ref{app:fields}, \ref{app:fields-certificates}, and~\ref{app:recourse},
respectively.
Appendix~\ref{app:qc} compares the results with convex quadratically
constrained problems, where the number of independent constraint Hessians
controls the arithmetic. Appendix~\ref{app:boundaries} records examples that
mark the limits of the methods: degenerate quartics, information-based
precision limits, penalty representations, and facial reduction.
Appendix~\ref{app:further-arithmetic} treats two settings outside the convex
classes above: quadratic constraints with few independent Hessians, and
cone constraints with bounded integer dimension over an explicitly
represented real number field. It develops common-field witness, finite-value,
attainment, and algebraic-input contracts, as summarized in
Section~\ref{sec:intro-further-arithmetic}.

Readers interested mainly in complexity can read Sections~\ref{sec:models},
\ref{sec:upper}, \ref{sec:reductions}, and~\ref{sec:constraints}; readers
interested in certificates, Sections~\ref{sec:models}
and~\ref{sec:algebraic}--\ref{sec:fields-certificates}; and readers
interested in solver design, Sections~\ref{sec:points}, \ref{sec:constraints},
\ref{sec:recourse}, and~\ref{sec:discussion}.
